% Level 1, Week 5 -- answer key
\answersweek{5}{Messages and My Blocks}{Answers for workbook chapter 5 and homework 5.}

\section{Lesson and workbook questions}
\begin{answer}[Exit ticket]
Nothing happens to the dog: \texttt{go} and \texttt{Go!} are two different messages, so the dog never hears its message. Fix: choose \texttt{go} in the dropdown of the dog's \enquote{when I receive} block.
\end{answer}

\begin{answer}[Reflection]
\enquote{polygon 6 60} tells you \textbf{what} it does (a hexagon with sides of 60) in one line; you do not have to read the repeat, move and turn blocks to understand it. Accessibility answers: chose colours from the colour-safe paint box, added a shape or word as well as colour, made text bigger, checked the project in grayscale.
\end{answer}

\section{Workbook exercises}
\begin{answer}[Exercise 5.1: Shape Machine]
\textbf{Expected behaviour:} each click on the flag erases the stage and draws four shapes that all start at the centre and hang below-right of it: a triangle (sides 120), a square (100), a hexagon (60) and an octagon (40). Adding \enquote{polygon 5 70} adds a pentagon.

\textbf{Question:} without the My Block, each shape needs pen down + repeat + move + turn + pen up = 5 blocks, so 4 shapes need about 20 blocks (accept 16--20). With the My Block the main script uses 4.

\textbf{What to look for:} the inputs are dragged from the define hat (not typed); \enquote{360 / sides} in the turn; the pink block is used at least 4 times with different numbers.
\end{answer}

\begin{answer}[Exercise 5.2: Stamp Garden]
\textbf{Expected behaviour:} each green-flag click erases the old garden, draws a thick bluish-green strip of grass along the bottom of the stage, then stamps 12 copies of the painted flower in random places, each a random size between 30\,\% and 100\,\% and each a different colour (the colour effect changes by 25 every stamp). The sprite itself then hides, so only the stamps are left. Every garden is different.

\textbf{What to look for:} two My Blocks (\texttt{flower} with an input, \texttt{grass} without); the main script reads like a plan; \enquote{show} at the start (otherwise the second garden has no flowers because the sprite is still hidden and stamping a hidden sprite stamps nothing).
\end{answer}

\begin{answer}[Exercise 5.3: Show Time]
\textbf{Expected behaviour:} on the flag the dancer hides and the presenter walks to the left and says \enquote{Welcome to the show!} for 2 seconds, then broadcasts \texttt{start show}. The dancer appears at (80, $-30$), changes costume 8 times (a short dance, about 2.4 seconds), then broadcasts \texttt{finale}, and the presenter says \enquote{Thank you for watching!}.

\textbf{Ada's bug:} the presenter speaks, then nothing happens; the dancer never appears. \textbf{Why:} the presenter broadcasts \texttt{start show} but the dancer is listening for a different message, \texttt{show start}, so no script receives \texttt{start show}. \textbf{Fix:} in the dancer's hat block choose \texttt{start show} from the dropdown.
\end{answer}

\section{Homework}
\begin{answer}[Core: Starry Sky]
\textbf{Expected behaviour:} on each flag click the stage is cleared, the cat hides, and 10 yellow five-pointed stars appear in random places, each with a random size between 10 and 60 steps. There are no lines joining the stars, because the pen is lifted before each \enquote{go to random position}. If there are joining lines, a \enquote{pen up} is missing (before the loop or at the end of the define script).

\textbf{Why 144 and not 72?} A five-pointed star goes all the way round \textbf{twice} while drawing, so the turns add up to 720 degrees: 720 ÷ 5 = 144. With 72 you get a pentagon, not a star.
\end{answer}

\begin{answer}[Challenge: moon]
One good plan: \enquote{when green flag clicked; erase all; hide; moon; repeat 10: go to random position, star (pick random 10 to 60)}. The \texttt{moon} define script: pen up, go to x: $-170$ y: 110, pen down, repeat 36 \{move 8, turn 10\}, pen up. Accept any moon position and size; look for a My Block with no input used once, and a main script that reads like a list of named parts.
\end{answer}

\begin{answer}[Creative: good night message]
Button: \enquote{when this sprite clicked $\rightarrow$ broadcast night falls}. Cat: \enquote{when I receive night falls $\rightarrow$ say Good night! for 2 seconds}. The two message names must be identical. A great answer also switches to a darker backdrop when the message arrives (a stage script \enquote{when I receive night falls $\rightarrow$ switch backdrop to \ldots}).
\end{answer}
